\section{Method: details}\label{sec:A04-method}

\subsection{Inexact kernel projection}\label{sec:A04-inexact}

\cref{prop:upper-bound} assumes that $\mathsf Z$ solves $\Kp\mathsf Z=\mathsf B$ exactly.
The CG solve of \cref{sec:04-proj} returns $\tilde{\mathsf Z}=\mathsf Z+\mathsf Z_{\rm err}$ with residual $\mathsf R=\Kp\tilde{\mathsf Z}-\mathsf B=\Kp\mathsf Z_{\rm err}$ and $\|\mathsf R_{:,j}\|_2\le\tau\|\mathsf B_{:,j}\|_2$ for every column.

\paragraph{The numerator is unaffected.}
$\Av=\Vmat^{\!\top}\Kmat\Vmat$ is evaluated directly, and $\Kmat\Gmat=0$.
The computed field is $\tilde{\Pmat}\Vmat=\Pmat\Vmat-\Gmat\mathsf Z_{\rm err}$, so its curl energy equals that of the exact projection: the residual gradient component $-\Gmat\mathsf Z_{\rm err}$ carries no curl energy.

\paragraph{The denominator is perturbed.}
The Ritz layer uses $\tilde{\mathsf M}_{\rm div}=\Vmat^{\!\top}\Mmat\Vmat-\mathsf B^{\!\top}\tilde{\mathsf Z}$.
Since $\mathsf B=\Kp\mathsf Z$, $\mathsf B^{\!\top}\mathsf Z_{\rm err}=\mathsf Z^{\!\top}\mathsf R$, hence
\[
  \tilde{\mathsf M}_{\rm div}=\Mdiv+\delta\mathsf M,\qquad \delta\mathsf M=-\mathrm{sym}(\mathsf Z^{\!\top}\mathsf R).
\]
For a coefficient vector $c$, $c^{\!\top}\delta\mathsf M\,c=-(\mathsf Zc)^{\!\top}\Kp(\mathsf Z_{\rm err}c)=-\langle\Gmat\mathsf Zc,\Gmat\mathsf Z_{\rm err}c\rangle_{\Mmat}$.
$\Gmat\mathsf Zc$ is the $\Mmat$-orthogonal projection of $\Vmat c$ onto $\mathrm{range}\,\Gmat$, so $\|\Gmat\mathsf Zc\|_{\Mmat}\le\|\Vmat c\|_{\Mmat}$, and $\|\Gmat\mathsf Z_{\rm err}c\|_{\Mmat}^2=\|\mathsf Rc\|_{\Kp^{-1}}^2$.
With $c^{\!\top}\Mdiv c=\|\Pmat\Vmat c\|_{\Mmat}^2=\mu(c)\|\Vmat c\|_{\Mmat}^2$, where $\mu(c)\in[0,1]$ is the divergence-free mass fraction of $\Vmat c$,
\[
  \eta(c):=\frac{|c^{\!\top}\delta\mathsf M\,c|}{c^{\!\top}\Mdiv c}\ \le\ \frac{\|\mathsf Rc\|_{\Kp^{-1}}}{\mu(c)\,\|\Vmat c\|_{\Mmat}}\ \le\ \frac{\lambda_{\min}(\Kp)^{-1/2}\,\|\mathsf Rc\|_2}{\mu(c)\,\|\Vmat c\|_{\Mmat}}.
\]

\paragraph{Consequence for the Ritz values.}
Let $\bar\eta=\sup\eta(c)<1$ over the coefficient subspace $\mathcal C$ kept by the Ritz layer.
Then $c^{\!\top}\tilde{\mathsf M}_{\rm div}c\le(1+\bar\eta)\,c^{\!\top}\Mdiv c$ on $\mathcal C$, and $\Pmat\Vmat$ is injective there.
Every computed Rayleigh quotient $c^{\!\top}\Av c/c^{\!\top}\tilde{\mathsf M}_{\rm div}c$ is therefore at least $(1+\bar\eta)^{-1}$ times the Rayleigh quotient of the exact field $\Pmat\Vmat c\in Z_h$.
The argument of Step~3 of the proof, applied to $k$-dimensional subspaces of $\mathcal C$ and their images under $\Pmat\Vmat$, gives
\[
  \tilde\theta_k\ \ge\ \lamh_k/(1+\bar\eta).
\]
The returned field $\tilde{\Pmat}\Vmat c$ has the same numerator and the larger mass $c^{\!\top}\Mdiv c+\|\mathsf Rc\|_{\Kp^{-1}}^2$, so its own Rayleigh quotient is lowered by a relative amount of second order in the residual.

\paragraph{How small is $\bar\eta$?}
The bound shows $\bar\eta\lesssim C\,\tau/\mu_{\min}$, where $\mu_{\min}\ge\tau_{\rm drop}$ is the smallest divergence-free fraction among the kept directions and $C$ collects the conversion between the Euclidean residual norm and the $\Kp^{-1}$ and $\Mmat$ norms, which depends on the mesh.
A priori, $\tau/\mu_{\min}$ can be as large as $10^{-2}$ before $C$ at $\tau=10^{-8}$ and $\tau_{\rm drop}=10^{-6}$, so $\bar\eta<1$ is not guaranteed.
In practice $\bar\eta$ can be computed after the fact: with $\mathsf Z$ from a direct solve of $\Kp$ it is the largest absolute eigenvalue of $\mathsf W^{\!\top}\mathsf D\,\delta\mathsf M\,\mathsf D\mathsf W$ on the kept Ritz subspace ($\mathsf W^{\!\top}\mathsf D\Mdiv\mathsf D\mathsf W=\mathsf I$), and \cref{sec:A04-stress} reports it for synthetic bases.
There $\bar\eta\approx(1\text{--}3)\,\tau$, i.e.\ the divergence-free fraction of the kept directions is far above $\tau_{\rm drop}$ except for the adversarial pure-gradient basis at $\tau=10^{-4}$.
Floating-point rounding is not covered by this argument.

\subsection{Rank check of assumption (A)}\label{sec:A04-rank}

\cref{tab:A04-rank} lists the check of (A) (\cref{eq:03-A}) for the operators of the model (free DOFs only, dense generalised eigenproblem): $\dim\ker\Kmat$ equals $\operatorname{rank}\Gmat$, equals $\operatorname{rank}\Kp$, equals the number of interior $H^1$ DOFs of order $p+1$ plus the extra wall potentials, in all $15$ cases.
The zero eigenvalues are at most $4\times10^{-16}$ of the largest, the first nonzero one at least $2.7\times10^{-4}$, so the threshold ($10^{-10}$) is not delicate; the kernel vectors lie in $\mathrm{range}\,\Gmat$ to $7\times10^{-14}$ (M-norm).
The smallest nonzero eigenvalue agrees with the label solver's to $10^{-13}$.
The meshes are small (at most $6312$ free DOFs) because the check is dense; the production meshes are not covered, except indirectly by $\Kmat\Gmat=0$ in the test suite.
% src: docs/35 §1 (command, setup, table, limitations)

\begin{table}[t]
\centering
\caption{Rank check of assumption (A). Meshes curved at geometry order $2$ (the pipeline uses order $3$); ``annulus'' is a solid torus ($b_1=1$), which adds no kernel in the E formulation, and the shell has two wall components (one extra potential). $\dim\ker\Kmat=\operatorname{rank}\Gmat=\operatorname{rank}\Kp$ in every row. ``Zero'': largest kernel eigenvalue relative to the largest eigenvalue; ``first'': smallest nonzero one; ``dist'': largest $\Mmat$-norm of the part of a kernel vector outside $\mathrm{range}\,\Gmat$.}
\label{tab:A04-rank}
\footnotesize
\setlength{\tabcolsep}{4pt}
\begin{tabular}{@{}lrrrrrccc@{}}
\toprule
mesh & tets & $p$ & free DOFs & $\dim\ker\Kmat$ & interior $+$ extra & zero & first & dist \\
\midrule
pillbox & 109 & 1 & 116 & 58 & $58+0$ & $3.4\!\times\!10^{-16}$ & $3.3\!\times\!10^{-2}$ & $1.7\!\times\!10^{-15}$ \\
pillbox & 109 & 2 & 672 & 282 & $282+0$ & $3.2\!\times\!10^{-16}$ & $7.5\!\times\!10^{-3}$ & $6.0\!\times\!10^{-15}$ \\
pillbox & 109 & 3 & 1996 & 781 & $781+0$ & $3.2\!\times\!10^{-16}$ & $2.9\!\times\!10^{-3}$ & $1.1\!\times\!10^{-14}$ \\
pillbox & 179 & 1 & 220 & 113 & $113+0$ & $3.6\!\times\!10^{-16}$ & $2.2\!\times\!10^{-2}$ & $2.4\!\times\!10^{-15}$ \\
pillbox & 179 & 2 & 1185 & 508 & $508+0$ & $3.3\!\times\!10^{-16}$ & $5.2\!\times\!10^{-3}$ & $6.3\!\times\!10^{-15}$ \\
pillbox & 179 & 3 & 3436 & 1367 & $1367+0$ & $3.4\!\times\!10^{-16}$ & $2.0\!\times\!10^{-3}$ & $1.4\!\times\!10^{-14}$ \\
shell & 139 & 1 & 184 & 93 & $92+1$ & $3.9\!\times\!10^{-16}$ & $2.5\!\times\!10^{-2}$ & $4.2\!\times\!10^{-15}$ \\
shell & 139 & 2 & 963 & 414 & $413+1$ & $3.3\!\times\!10^{-16}$ & $5.7\!\times\!10^{-3}$ & $8.6\!\times\!10^{-15}$ \\
shell & 139 & 3 & 2756 & 1103 & $1102+1$ & $3.2\!\times\!10^{-16}$ & $2.1\!\times\!10^{-3}$ & $3.9\!\times\!10^{-14}$ \\
annulus & 78 & 1 & 76 & 38 & $38+0$ & $3.5\!\times\!10^{-16}$ & $7.7\!\times\!10^{-2}$ & $1.2\!\times\!10^{-15}$ \\
annulus & 78 & 2 & 462 & 192 & $192+0$ & $3.0\!\times\!10^{-16}$ & $1.3\!\times\!10^{-2}$ & $4.4\!\times\!10^{-15}$ \\
annulus & 78 & 3 & 1392 & 540 & $540+0$ & $2.2\!\times\!10^{-16}$ & $3.9\!\times\!10^{-3}$ & $1.1\!\times\!10^{-14}$ \\
pillbox with pipes & 316 & 1 & 436 & 229 & $229+0$ & $2.4\!\times\!10^{-16}$ & $4.2\!\times\!10^{-3}$ & $7.3\!\times\!10^{-15}$ \\
pillbox with pipes & 316 & 2 & 2220 & 969 & $969+0$ & $1.4\!\times\!10^{-16}$ & $7.7\!\times\!10^{-4}$ & $3.1\!\times\!10^{-14}$ \\
pillbox with pipes & 316 & 3 & 6312 & 2547 & $2547+0$ & $1.4\!\times\!10^{-16}$ & $2.7\!\times\!10^{-4}$ & $7.2\!\times\!10^{-14}$ \\
\bottomrule
\end{tabular}
\end{table}

\subsection{Stress test of the bound under the inexact projection}\label{sec:A04-stress}

The Ritz layer of \cref{sec:04-ritz} was fed synthetic bases on a stored label geometry whose exact discrete eigenpairs are known (test geometry 24: $4818$ tetrahedra, $51\,105$ DOFs, $18\,662$ potentials; geometry 27: $6315$ tetrahedra, $69\,645$ DOFs), with $m=12$ columns, $K=4$ and the configuration-file defaults except the CG tolerance $\tau$.
The bases are: (a) the four exact eigenvectors plus eight random columns; (b) the same with relative noise $\epsilon\in\{0,10^{-8},10^{-6},10^{-4}\}$ on the eigenvectors; (c) twelve random columns; (d) the exact eigenvectors plus a gradient $\Gmat r$ of mass fraction $\gamma=\|g\|^2/(\|u\|^2+\|g\|^2)$, plus eight random columns; (g) the four eigenvectors plus eight pure-gradient columns.
We report $\min_k(\theta_k/\lamh_k-1)$ over four modes and three seeds, the number of violations ($<-10^{-12}$) out of $12$, the a posteriori $\bar\eta$ of \cref{sec:A04-inexact}, the number of discarded directions, and the number of spurious Ritz values ($\theta<0.5\,\lamh_1$) of $\spn(\Vmat)$ without the projection.
The stored eigenvectors are float32, i.e.\ already ``exact plus $6\times10^{-8}$ noise'', so margins at $\epsilon=0$ are $\pm10^{-14}$ rather than $0$.

\begin{table}[t]
\centering
\caption{Stress test, test geometry 24 (three seeds, four modes): smallest margin $\theta_k/\lamh_k-1$ and, after the slash, the number of violations out of $12$, for the CG tolerance $\tau$; the shipped value is $10^{-8}$. Bases (a) and (d) with $\gamma\le10^{-6}$ are at rounding level for every $\tau$ and omitted; (b) with noise $\epsilon\le10^{-6}$ is at rounding level or above, with $\epsilon=10^{-4}$ it gives $+10^{-6}$, and (c) gives $+2.8\times10^{2}$. ``Spurious'': Ritz values of $\spn(\Vmat)$ below $0.5\,\lamh_1$ without the projection (with the projection the count is $0$ for $\tau\le10^{-6}$).}
\label{tab:A04-stress}
\scriptsize
\setlength{\tabcolsep}{3pt}
\begin{tabular}{@{}lccccc@{}}
\toprule
basis & $\tau=10^{-4}$ & $10^{-6}$ & $10^{-8}$ & $10^{-10}$ & spurious \\
\midrule
(a) exact $+$ random & $-5\!\times\!10^{-15}$/0 & $-2\!\times\!10^{-14}$/0 & $-2\!\times\!10^{-15}$/0 & $-4\!\times\!10^{-14}$/0 & 0 \\
(d) $\gamma=10^{-4}$ & $-7.3\!\times\!10^{-11}$/3 & $-7.7\!\times\!10^{-13}$/0 & $-4.5\!\times\!10^{-14}$/0 & $-3.8\!\times\!10^{-13}$/0 & 0 \\
(d) $\gamma=10^{-2}$ & $-7.7\!\times\!10^{-9}$/8 & $-6.6\!\times\!10^{-11}$/4 & $-1.3\!\times\!10^{-12}$/1 & $-2.7\!\times\!10^{-13}$/0 & 0 \\
(d) $\gamma=0.5$ & $-7.5\!\times\!10^{-7}$/7 & $-2.6\!\times\!10^{-9}$/6 & $-2.0\!\times\!10^{-11}$/4 & $-6.1\!\times\!10^{-13}$/0 & 1 \\
(d) $\gamma=0.9$ & $-6.2\!\times\!10^{-6}$/4 & $-2.6\!\times\!10^{-8}$/6 & $-6.1\!\times\!10^{-10}$/5 & $-5.7\!\times\!10^{-12}$/2 & 4 \\
(d) $\gamma=0.99$ & $-7.4\!\times\!10^{-5}$/4 & $-4.5\!\times\!10^{-7}$/3 & $-6.6\!\times\!10^{-9}$/6 & $-1.4\!\times\!10^{-10}$/5 & 4 \\
(g) pure gradient & $-1.0$/12 & $+3\!\times\!10^{-14}$/0 & $+3\!\times\!10^{-14}$/0 & $+3\!\times\!10^{-14}$/0 & 8 \\
\bottomrule
\end{tabular}
\end{table}

\begin{table}[t]
\centering
\caption{A posteriori perturbation $\bar\eta$ and CG iterations of the stress test. Maximum over bases (a)--(d) for geometry 24 (geometry 27 in parentheses), and $\bar\eta$ for the pure-gradient basis (g).}
\label{tab:A04-eta}
\scriptsize
\setlength{\tabcolsep}{3pt}
\begin{tabular}{@{}lcccc@{}}
\toprule
 & $\tau=10^{-4}$ & $10^{-6}$ & $10^{-8}$ & $10^{-10}$ \\
\midrule
max $\bar\eta$, (a)--(d) & $2.5\!\times\!10^{-4}$ ($3.7\!\times\!10^{-4}$) & $2.5\!\times\!10^{-6}$ ($2.0\!\times\!10^{-6}$) & $2.6\!\times\!10^{-8}$ ($2.8\!\times\!10^{-8}$) & $2.8\!\times\!10^{-10}$ ($4.3\!\times\!10^{-10}$) \\
$\bar\eta$, (g) & $1.0$ & $3\!\times\!10^{-24}$ & $4\!\times\!10^{-25}$ & $4\!\times\!10^{-25}$ \\
max CG iterations & 80 (50) & 120 (90) & 160 (120) & 200 (150) \\
\bottomrule
\end{tabular}
\end{table}

\paragraph{Reading.}
(i) With an exact projection (sparse direct $\Kp^{-1}$) the smallest margin over all bases except random is $-2.3\times10^{-13}$ (geometry 24) and $-2.9\times10^{-13}$ (geometry 27): no violation at the $10^{-12}$ level.
(ii) With the inexact CG projection the plain bound is violated only for bases with a gradient admixture, and the violation grows with $\gamma$ and $\tau$; at the shipped $\tau$ it is at most $1.5\times10^{-8}$ ($\gamma=0.99$), $2\times10^{-11}$--$1.4\times10^{-10}$ at $\gamma=0.5$ and $1.3\times10^{-12}$ at $\gamma=10^{-2}$ (geometry 24; the ranges are over both geometries), and at $\tau=10^{-4}$ up to $8.0\times10^{-5}$.
A trained basis has a gradient mass fraction of about $0.5$--$0.6$ before projection (the lowest-order variant measured $0.59$), so the $\gamma=0.5$ row is the relevant one.
(iii) The violation never exceeds $\bar\eta$, and the a posteriori guarantee $\tilde\theta_k\ge\lamh_k/(1+\bar\eta)$ held in all $156+104$ runs (geometries 24 and 27) with $\bar\eta<1$, i.e.\ all except the $3+2$ runs of the pure-gradient basis at $\tau=10^{-4}$.
(iv) The discarding rule works when the residual gradient left by CG has mass fraction $(C\tau)^2<\tau_{\rm drop}=10^{-6}$: at $\tau=10^{-4}$ only $4$--$5$ of the $8$ pure-gradient directions are discarded, the others enter the Ritz problem with $\theta\approx0$ ($\bar\eta=1$, margin $-100\,\%$, guarantee void), whereas at $\tau\le10^{-6}$ all $8$ are discarded.
(v) Without the projection, eight pure-gradient columns give eight Ritz values at $\theta\approx0$ and the admixtures $\gamma=0.9,0.99$ give four spurious values; with the projection the count is $0$ for $\tau\le10^{-6}$.
White-noise bases (random padding, basis (c)) do not produce spurious values even without the projection, because their curl energy is large: the projection matters for smooth, gradient-rich bases.
% src: docs/35 §2 (command, setup, both tables, further facts i-vi, reading); scripts/paper/check_ritz_bound.py
\paragraph{Limitations.}
Two geometries of one family, one preconditioner (Jacobi), the iteration cap $3000$ never reached; the bases are synthetic (trained bases have a continuum of $\mu$ values, and the pure-gradient and exact-plus-gradient rows are adversarial extremes); $\bar\eta$ is evaluated over the whole kept subspace, which is conservative; mode identity is not tested.

\subsection{Cost and default configuration}\label{sec:A04-cost}

\paragraph{Default configuration.}
The configuration-file defaults of the model are embedding dimension $d=128$, $L=4$ attention blocks, $4$ heads, basis width $m=24$, $K=6$ modes and a $64$-dimensional Fourier map of length scale $0.1$.
% src: configs/field_3d.yaml (embed_dim 128, n_layers 4, n_heads 4, n_basis 24, num_field_modes 6, rff_dim 64, rff_length_scale 0.1)
The reported experiments use a smaller model (\cref{sec:experiments}).

\paragraph{Operator sizes.}
For a curved pillbox-with-pipes mesh of $10.5$k tetrahedra, order two has $112$k DOFs, $41$k potentials and $148$k tokens, and order three $282$k DOFs, $100$k potentials and $253$k tokens; assembling the operators takes $7$--$10$\,s and $37$\,s on one thread.
% src: docs/28 §4 table (netgen curvaturesafety 2: 10.5k tets, p2 112k DOF / 41k potentials / 148k tokens / 7-10 s; p3 282k / 100k / 253k / 37 s)
Operators are rebuilt per sample rather than stored, because they would take about $0.4$\,GB per sample at order two against $1.4$\,MB for the curved mesh.
% src: docs/28 §4 Storage (1.4 MB per sample; ~0.4 GB per sample at p2)

\paragraph{Training step.}
For a mesh of $13$k tetrahedra ($183$k tokens, order two) on four CPU threads we measured $4.2$\,s for operator assembly, $0.34$\,s for batch collation and about $22$\,s per training step.
% src: docs/28 §4 "Measured per step" (13k tets, 183k tokens, CPU, 4 threads)
Gradient checkpointing (recomputing the encoder and each attention and feed-forward half in the backward pass) lowered the peak process memory of training from $9.9$ to $3.8$\,GB with bit-identical gradients.
% src: docs/28 §4 (training memory: 9.9 GB -> 3.8 GB, gradients bit-identical)

\paragraph{CG iterations.}
A CG iteration costs $O(\mathrm{nnz}(\Kp)\,m)$ in float64 and the backward pass adds one more solve; the $m\times m$ Ritz layer is negligible.
The iteration count grows roughly like $h^{-1}$: in the lowest-order variant it rose from $60$ to $170$ while the edge count grew from $2\,053$ to $47\,919$.
% src: docs/20 §5 (PCG iterations 60 ... 170 for N_e 2053 ... 47919)
